% Independent distribution audit, 2026-10-08. No frozen source changed.
\subsection{Actual finite atoms: distribution with a moving target}
\label{sec:actual-k-distribution-joint-audit}

This is a fresh hypothesis audit of the distribution step, for the
current intended coefficient $K=67619/10000$. It does not reprove the
external dispersion theorem, the finite-tree density, or the endpoint
factorization lemma. The conclusion below is an implication from
explicit finite-atom data, rather than from a saved numerical status.
We found no counterexample to that implication. We supply the missing
divisor-weighted justification of square-collision removal explicitly.

\paragraph{Primary statements checked anew.}
On 8 October 2026 the primary arXiv record for
\href{https://arxiv.org/abs/2309.08522}{Lichtman 2309.08522}
still exposes only v1. We read its Definition 2.1, Proposition 5.2,
Proposition 3.2, and Lemma 3.4, including their proofs where needed.
Fresh HTML and plain-text copies are saved as
\texttt{external-source-audit/l23-v1-joint-current.*}.
The relevant text starts at extracted lines 1239, 2411, 1316, and 1370.
We separately read
\href{https://arxiv.org/html/2006.07088v1}{Maynard II v1},
Lemma 7.1 and the proof of Proposition 5.2; the saved text is
\texttt{maynard-II-v1-joint-current.txt}, lines 1283 and 1386--1467.
In particular, lines 1414--1423 explicitly permit arbitrary intervals
in both long variables. Its underlying Lemma 7.1 takes a supremum over
all primitive residues, so this extension does not import the
fixed-residue restriction of other Maynard theorems.

For low-rank atoms we also checked the generalized BV statement in
\href{https://arxiv.org/html/1407.4897v4}{Polymath 1407.4897v4},
Claim 2.6 and Theorem 2.8 (saved text lines 1081 and 1265).
These correspond to Claim 12 and Theorem 14 in the journal numbering.
Only its unconditional range below $1/2$ is used. L23's arXiv theorem
is an external analytic input; this audit does not relabel it as a
verified journal publication or independently verify its spectral proof.

\paragraph{The precise atom input.}
Fix a positive outer cutoff $a$, a maximum number $k_*$ of original
prime slots, and one finite buffered modulus strategy. Write $w=N^\omega$
with fixed $\omega>0$. An atom has the form
\begin{equation}\label{eq:actual-k-atom-input}
 c_J\,\mathbf1_{\rm squarefree}
       (\chi_{J_1}*\cdots*\chi_{J_r})*\beta_J,
 \qquad r\le M_*,\quad |\beta_J|\le1,
\end{equation}
with prime intervals above $w$, small-prime support for $\beta_J$, and
the appropriate fixed parity choice of the latter. The finite sum of
$|c_J|$ is independent of $N$. All prime divisors of every modulus
are coprime to $N$. The repeated-label factorial is already included
in $c_J$, once. The complete support is at most $N^{\theta_*}$ for
some fixed $\theta_*<3/5$; atoms of rough rank at most two have support
at most $N^{2(19/96)+o(1)}<N^{1/2-c}$.
After the finite strategy is fixed, choose $\kappa$ also so small that
$\theta_* /(1-\kappa)<3/5$ and that the complete low-rank exponent
divided by $1-\kappa$ is below $1/2$. Thus the asserted modulus-level
bounds remain valid in the local $x$ coordinate, not only in base $N$.

For every local product scale $x\asymp\prod p_i\ge N^{1-\kappa}/O(1)$
and every medium length $U$, the finite-tree/factorization step must
supply a partition of the \emph{labelled slots}, fixed on the entire
box. Let $B_1,B_2,B_3$ be upper bounds for the three resulting rough
products, and put $A_0=Ux^{-\zeta}$. Before fine-box and beta losses,
the required output is
\begin{align}\label{eq:actual-k-factorization-input}
 B_1&\le A_0,&
 A_0^2B_2B_3^2&\le x^{1-3\zeta},\nonumber\\
 N^\sigma A_0B_1B_2^5B_3^2&\le x^{2-3\zeta},&
 A_0^{1+2\sigma}B_1B_2^{5-2\sigma}B_3^{2+2\sigma}
 &\le x^{2-2\zeta},\quad \sigma=7/32.
\end{align}
These statements are stronger and more specific than pointwise
membership of each modulus in a well-factorable support set. Their
derivation for the actual fourth/fifth/sixth strategy is the separate
tree/factorization audit; no arbitrary restriction of a factorable
sequence is presumed factorable here.

For a group with $r_i$ prime slots, divide its unrestricted convolution
by $r_i!$. Each resulting coefficient is one-bounded, since an integer
has at most $r_i!$ ordered representations in those slots, including
when primes repeat. Multiply the atom scalar by $r_1!r_2!r_3!$.
This restores exactly the original atom; it does not insert a second
repeated-label factorial. Append $\beta_J$ to the third group. The
small/rough factorization of each integer is unique, so this group
remains one-bounded. Empty groups have coefficient supported at $1$.
Usual dyadic decompositions retain that singleton explicitly; there
is no use of a Siegel--Walfisz claim for a unit slot.

\paragraph{A repair that is valid after Heath--Brown expansion.}
The elementary count $N/d+1$ in the older collision display is a
bound for unweighted integers. It is not literally a bound for an
HB convolution carrying a divisor weight. The following argument
supplies exactly the required stronger justification without a new
distribution theorem. Let $D\le N^{\theta_*}$, and let
\[
 \mathcal C=\{d\le D:p^2\mid d\text{ for some prime }p>w\}.
\]
For any fixed divisor order, both the total HB coefficient at an
integer $m\ll N$ and the absolute modulus coefficient are $N^{o(1)}$.
The same is true of $d/\varphi(d)$. Moreover
\[
 \sum_{d\in\mathcal C}\frac1d\ll\frac{\log N}{w},
 \qquad \#\mathcal C\ll\frac D w,
\]
by summing over $d=p^2e$ and allowing multiple counting. Thus the
\emph{whole} progression and exact-coprime mean errors from these
collisions satisfy
\begin{equation}\label{eq:actual-k-hb-collision-repair}
 E_{\rm coll}\ll N^{o(1)}
      \sum_{d\in\mathcal C}\left(\frac Nd+1+\frac N{\varphi(d)}\right)
 \ll\frac{(N+D)N^{o(1)}}w
 \ll N^{1-\omega/2}.
\end{equation}
The last inequality holds for all sufficiently large $N$ because
$\omega$ and the divisor orders are fixed. This is more than enough
for every inverse logarithmic saving. It also covers $m=N$; for that
integer the congruence holds for every modulus, and the $+1$ term
was retained. No claim of uniformity as $\omega\downarrow0$ is made.
The repair loses no limiting sieve constant.

\paragraph{Heath--Brown retains the original groups.}
Work first with a smooth outer prime-product weight, of fixed
polylogarithmic normalized derivative bounds. Apply the order-three
HB identity separately at the original prime scale $R_i$.
Discard a dyadic HB rectangle unless its original group meets that
prime interval. Its at most six dyadic lower endpoints then have
product at least $R_i/2^6\gg N^a$. This persists after Mellin
separation, even outside the original curved support.

Keep the original prime roughness condition
$P^-(\prod p_i)>z_0=x^{1/(\log\log x)^3}$ before expansion. It is
automatic for the primes being counted, and makes every nonunit
HB factor $z_0$-rough. Every Mobius slot has scale
$\ll R_i^{1/3}\ll x^{1/3}$, with fixed multiplicative constants.
Consequently the following dichotomy is exhaustive for large $N$:
\begin{enumerate}
\item An individual HB slot lies in
$[x^{2\zeta},x^{1/3+\zeta/2}]$. Its coefficient is rough $1$ or
rough $\mu$, with an interval, a normalized logarithm, and a Mellin
twist. It has the full auxiliary-coprimality SW property proved in
\texttt{focused-rough-sw-product-interface-audit.tex}. No grouped
convolution is substituted for this individual SW slot.
\item Every slot is short, below $x^{2\zeta}$, or long, above
$x^{1/3+\zeta/2}$. Each original group has a long slot: otherwise
its product is at most $x^{12\zeta+o(1)}<N^a/O(1)$ when
$12\zeta<a/2$. Three long slots would have product exceeding
$x^{1+3\zeta/2}$, whereas the whole dyadic product is $O_{k_*}(x)$.
There are therefore at most two original prime groups.
\end{enumerate}
This rules out a hidden residual product $x^{O(k_*\zeta)}$ in the
exceptional case: that case itself forces $k\le2$.

\paragraph{Every Proposition 5.2 hypothesis.}
Group the other HB factors into $\beta_m$. Its coefficient has fixed
divisor order. Truncate only this complementary coefficient at a
sufficiently large logarithmic threshold, then normalize it; the
selected SW coefficient need not be truncated. The discarded actual
progression contribution is bounded, with arbitrary logarithmic
saving, by
\[
 (\log N)^{-B}\sum_{v\ll N,\ v\ne N}
       \tau_C(v)\tau_C(|N-v|)\ll N(\log N)^{C'-B}.
\]
Cauchy--Schwarz and fixed divisor moments prove the last estimate.
Here $v$ is the full HB product, not the complementary factor $m$.
The isolated $v=N$ contribution is $N^{\theta_*+o(1)}$; the mean
uses $\sum_{d\le D}\tau_C(d)/\varphi(d)\ll\log^{O_C(1)}N$.

Normalize the selected logarithm if present. The resulting
$\alpha_n,\beta_m$ are one-bounded, $UV\asymp x$, and the individual
$\alpha_n$ has precisely L23 Definition 2.1, including arbitrary
auxiliary coprimality and subintervals. The three modulus coefficients
are one-bounded by the preceding atom argument and supported on
$(q_i,N)=1$. They may depend on $N$: the theorem is uniform for
these coefficient classes, with the SW constants just verified
uniformly.

Substituting $U=A_0x^\zeta$ into
\eqref{eq:actual-k-factorization-input} leaves savings
$\zeta$, $2\zeta$, and $(1-2\sigma)\zeta=9\zeta/16$.
After the fixed fine-box, small-factor and local-scale losses choose
$\epsilon_{\rm disp}$ much smaller than these savings. The four
conditions in the \emph{actual} source become
\begin{align*}
 Q_1&<Ux^{-\epsilon_{\rm disp}},&
 U^2Q_2Q_3^2&<x^{1-8\epsilon_{\rm disp}},\\
 N^\sigma UQ_1Q_2^5Q_3^2&<x^{2-15\epsilon_{\rm disp}},&
 U^{1+2\sigma}Q_1Q_2^{5-2\sigma}Q_3^{2+2\sigma}
 &<x^{2-15\epsilon_{\rm disp}}.
\end{align*}
All factor powers in the final expression are positive; in particular
$5-2\sigma=73/16>0$. Passing from support endpoints to any dyadic
subbox therefore preserves the inequalities. Also
$A_0\in[x^\zeta,x^{1/3-\zeta/2}]$, exactly the factorization range.
For $\kappa<\epsilon_{\rm disp}/(2(1+\epsilon_{\rm disp}))$ one has
$N<x^{1+\epsilon_{\rm disp}}$ on every retained local scale, eventually.
The explicit residue supremum in L23 Proposition 5.2 thus contains
the actual target $N$. There is no exceptional set of even targets.

\paragraph{One long slot.}
There is only one original group. The other at most five factors have
product $M\le x^{10\zeta+o(1)}$, so the long scale is
$T\gg x^{1-10\zeta-o(1)}$. Choose the parameters so that
$1-10\zeta>3/5+2\epsilon_{\rm disp}$. For each remaining product
$m$, terms with $(m,d)>1$ vanish in both the progression and its
primitive mean, since $(N,d)=1$. Otherwise use residue
$Nm^{-1}\pmod d$. The inequality $T\ge d x^{\epsilon_{\rm disp}}$
allows Lemma 3.4 at both endpoints of any subinterval of $[T,2T]$.

Its displayed unrestricted rough mean can be replaced by the exact
coprime mean: the difference is at most
\[
 \frac1{\varphi(d)}\sum_{p\mid d,\ p>z_0}(2T/p+1)
 \ll_A \frac{T}{d\log^A x}.
\]
Indeed $d/\varphi(d)\ll\log\log(3d)$, $\omega(d)\ll\log x$,
$z_0$ exceeds every fixed logarithmic power, and $d/T\le x^{-\epsilon}$.
Summing this error with the fixed-order residual coefficients and
the modulus coefficients loses only logarithmic powers. A smooth
logarithm or twist is handled by partial summation. Modulus $1$ has
zero discrepancy separately.

\paragraph{Two long slots and arbitrary subintervals.}
Their lengths $T_1,T_2$ exceed $x^{1/3+\zeta/2}$. If there are two
original groups, the at most ten residual variables have product
$M\le x^{20\zeta+o(1)}$; for one original group the stronger
$M\le x^{8\zeta+o(1)}$ holds. Take parameters so small that
\[
 3/5<(1-20\zeta)(2/3-3\epsilon_{\rm disp}),
 \qquad 1/3+\zeta/2>3\epsilon_{\rm disp}.
\]
Then L23 Proposition 3.2 applies to the entire modulus range and has
uniformity over every nonzero target residue. Its absolute modulus
sum permits the original $|\rho_N(d)|\le1$ without an atom
factorization; alternatively each fixed atom has bounded total
coefficient and may be treated separately. The residual $m$ coefficient
is divisor-bounded, as required. The two long factors are rough
integers with coefficient $1$, before smooth partial summation.

The needed interval extension is explicit in Maynard's proof quoted
above. Consequently two partial summations allow separate Mellin
twists and logarithms on both long slots, at a cost bounded by the
product of their total variations, hence a fixed logarithmic power.
This is \emph{not} a claim that arbitrary arithmetic coefficients on
the two long slots are allowed. The source proof uses smooth functions
with polylogarithmic derivative bounds and a pointwise primitive-
residue estimate; it therefore preserves the required uniformity.

\paragraph{Low ranks, tails, and prime powers.}
For rank at most two, generalized BV applies below $x^{1/2-c'}$.
For $k=1$ use classical BV. For $k\ge2$ select one original prime
slot and combine the rest. Both lengths are between fixed positive
powers of $x$ and their complementary powers. The prime slot,
restricted to any interval and with a smooth twist, satisfies the
auxiliary SW condition by ordinary prime SW and deletion of primes
dividing the auxiliary integer. The other coefficient is
divisor-bounded. Complex weights follow by four real/imaginary
applications of the real-valued statement. There is no use of GEH
beyond $1/2$.

The finite number of atom labels and HB patterns, and the
polylogarithmic number of dyadic rectangles, are absorbed by starting
with a larger logarithmic saving. Smooth Mellin tails are bounded
after the modulus sum by the same shifted-divisor moment above.
Proper prime powers in an original slot contribute
$N^{1-a/2+o(1)}$ before and after the divisor bound for $m\ne N$;
the isolated equality $m=N$ costs only $N^{\theta_*+o(1)}$.
Thus removal of prime powers does not multiply their total mass by
the modulus level. Partial summation removes the prime logarithms.

\paragraph{The closed distribution implication.}
For the \emph{same} compiled coefficients and every permitted smooth
outer weight, the preceding argument proves
\begin{equation}\label{eq:actual-k-joint-signed-distribution}
 \sum_d\rho_N(d)\left\{
  \sum_{\prod p_i\equiv N\ (d)}W_N(\boldsymbol p)
  -\frac1{\varphi(d)}\sum_{(\prod p_i,d)=1}W_N(\boldsymbol p)
 \right\}\ll_A N\log^{-A}N.
\end{equation}
All constants may depend on the fixed cutoff, finite strategy,
smoothing precision and number of slots; they do not depend on
the even target. The proof does not impose a new numerical loss in
$K$. An atom assignment may be chosen separately for a Type II scale,
but only as an exact representation of the already fixed weight.
It may not be selected from the observed signs of progression errors.

\paragraph{Cross-check of the positive assembly.}
We additionally read \texttt{actual-k-positive-assembly.tex}. Its
coprime-mass union bound, lines 46--68, correctly uses
$\sum_{d:p\mid d}|\rho_N(d)|/\varphi(d)\ll\log N/(p-1)$,
with no factor $N^{\theta_*}$. Its pointwise upper sum is used before
distribution, including at $N-rq=0$; this avoids a signed-boundary
argument. The omitted primes $N-rq<N^\rho$ contribute $O(N^\rho)$
for ordered prime pairs. The added low-product paragraph, lines
84--101, also correctly covers the part of the global radial majorant
with $m<2N^{1-\kappa}$: its absolute progression sum is at most
$\sum_{m<2N^{1-\kappa}}O_W(\tau_C(m)\tau(N-m))
 =N^{1-\kappa+o(1)}$, since $m=N$ cannot occur there eventually.
The exact-coprime mean loses at most another logarithmic factor.
A smooth cutoff $\chi(m/N^{1-\kappa})$ has bounded multiplicative
derivatives on every dyadic chart, so this removal does not require
distribution below the proved local-scale range. For its bad domain, putting
$\beta=b/(1+b)$ gives $r>(rq)^\beta$ and $q\ge\sqrt{rq}$,
so the fixed cutoff is available. The radial limiting density is
$dt\,dy/[t(1-t)]$ for $y=rq/N$, and the $y<\eta$ tail is
$O_a(\eta N/\log N)$. Hence its integral is $\log(1/b)$.
These steps correctly yield the bad-count implication once the
same-tree majorant and density inputs are supplied.

\paragraph{Scope of the result.}
The audit closes the HB, moving-residue, long-variable, subinterval,
and divisor-weighted collision interfaces of the finite-atom
distribution argument. It does not certify the numerical value of
$K_{\rm end}$, nor independently prove that the actual prefix tree
supplies \eqref{eq:actual-k-atom-input} and
\eqref{eq:actual-k-factorization-input}; those are explicitly separate
joint-audit tasks. No new distribution hypothesis is required after
those data are supplied. Arbitrary arithmetic outer restrictions,
cutoffs tending to zero with $N$, or a number of slots growing with
$N$ are outside the proved implication. In particular, selecting only
one primitive residue class modulo $3$ in an outer prime slot would
destroy its SW property and is not licensed by bounded multiplicity.
